
\documentclass{article}
\usepackage{colortbl}
\usepackage{makecell}
\usepackage{multirow}
\usepackage{supertabular}

\begin{document}

\newcounter{utterance}

\centering \large Interaction Transcript for game `codenames', experiment `risk\_high', episode 1 with glm{-}5.2.
\vspace{24pt}

{ \footnotesize  \setcounter{utterance}{1}
\setlength{\tabcolsep}{0pt}
\begin{supertabular}{c@{$\;$}|p{.15\linewidth}@{}p{.15\linewidth}p{.15\linewidth}p{.15\linewidth}p{.15\linewidth}p{.15\linewidth}}
    \# & $\;$A & \multicolumn{4}{c}{Game Master} & $\;\:$B\\
    \hline

    \theutterance \stepcounter{utterance}  
    & & & \multicolumn{2}{p{0.3\linewidth}}{
        \cellcolor[rgb]{0.9,0.9,0.9}{
            \makecell[{{p{\linewidth}}}]{
                \texttt{\tiny{[GM$|$GM]}}
                {'hidden': {'team': ['verificação', 'piano', 'pato', 'gelo', 'ópera', 'morcego', 'passe', 'coroa', 'tronco'], 'innocent': ['tela', 'escorpião', 'mina'], 'opponent': ['gancho', 'maçã', 'polvo', 'ponto', 'romance', 'urso', 'verme', 'laboratório'], 'assassin': ['estádio', 'centro', 'bola', 'igreja', 'piscina']}, 'revealed': {'team': {'team': [], 'innocent': [], 'opponent': [], 'assassin': []}, 'opponent': {'team': [], 'innocent': [], 'opponent': [], 'assassin': []}}}
            }
        }
    }
    & & \\ \\

    \theutterance \stepcounter{utterance}  
    & & \multicolumn{4}{p{0.6\linewidth}}{
        \cellcolor[rgb]{0.9,0.9,0.9}{
            \makecell[{{p{\linewidth}}}]{
                \texttt{\tiny{[P1$\langle$GM]}}
                \texttt{Forneça uma pista de uma única palavra que se relacione com algumas das palavras da sua equipa da lista abaixo. Pode escolher quantas e quais palavras a sua pista relaciona.} \\
\texttt{O objetivo é fazer com que o seu parceiro adivinhe todas as palavras da sua equipa com base nas pistas que lhe dá, antes que a outra equipa tenha descoberto todas as suas. Se o seu parceiro adivinhar uma palavra 'oponente' ou 'distrator', a vez de adivinhar termina e a equipa adversária passa a jogar. Se o seu parceiro adivinhar uma palavra 'matar', o jogo termina e a sua equipa perde imediatamente.} \\
\texttt{Quanto mais palavras a sua pista relacionar, mais rapidamente poderá ganhar o jogo, mas mais difíceis serão as associações para o seu parceiro adivinhar, por isso escolha a sua palavra{-}pista com sabedoria. A palavra{-}pista tem de estar semanticamente relacionada com as palavras{-}alvo, não pode ser uma das palavras das listas nem conter partes delas.} \\
\\ 
\texttt{Forneça sempre a sua pista de uma única palavra e a sua lista de palavras{-}alvo relacionadas separadas por vírgulas no seguinte formato e torne as suas respostas tão curtas quanto possível, nunca inclua qualquer outro texto além do exigido neste formato:} \\
\\ 
\texttt{DICA: <PALAVRA>} \\
\texttt{ALVOS: <ALVOS>} \\
\\ 
\texttt{As palavras da sua equipa são: verificação, piano, pato, gelo, ópera, morcego, passe, coroa, tronco.} \\
\texttt{As palavras da equipa adversária são: gancho, maçã, polvo, ponto, romance, urso, verme, laboratório.} \\
\texttt{As palavras distratoras são: tela, escorpião, mina.} \\
\texttt{As palavras matar são: estádio, centro, bola, igreja, piscina.} \\
            }
        }
    }
    & \\ \\

    \theutterance \stepcounter{utterance}  
    & \multicolumn{4}{p{0.6\linewidth}}{
        \cellcolor[rgb]{0.9,0.9,0.9}{
            \makecell[{{p{\linewidth}}}]{
                \texttt{\tiny{[P1$\rangle$GM]}}
                \texttt{DICA: música} \\
\texttt{ALVOS: piano, ópera} \\
            }
        }
    }
    & & \\ \\

    \theutterance \stepcounter{utterance}  
    & & & \multicolumn{2}{p{0.3\linewidth}}{
        \cellcolor[rgb]{0.9,0.9,0.9}{
            \makecell[{{p{\linewidth}}}]{
                \texttt{\tiny{[GM$|$GM]}}
                \texttt{música} \\
            }
        }
    }
    & & \\ \\

    \theutterance \stepcounter{utterance}  
    & & & \multicolumn{2}{p{0.3\linewidth}}{
        \cellcolor[rgb]{0.9,0.9,0.9}{
            \makecell[{{p{\linewidth}}}]{
                \texttt{\tiny{[GM$|$GM]}}
                ['piano', 'ópera']
            }
        }
    }
    & & \\ \\

    \theutterance \stepcounter{utterance}  
    & & & \multicolumn{2}{p{0.3\linewidth}}{
        \cellcolor[rgb]{0.9,0.9,0.9}{
            \makecell[{{p{\linewidth}}}]{
                \texttt{\tiny{[GM$|$GM]}}
                {'word': 'piano', 'assignment': 'team'}
            }
        }
    }
    & & \\ \\

    \theutterance \stepcounter{utterance}  
    & & & \multicolumn{2}{p{0.3\linewidth}}{
        \cellcolor[rgb]{0.9,0.9,0.9}{
            \makecell[{{p{\linewidth}}}]{
                \texttt{\tiny{[GM$|$GM]}}
                {'word': 'ópera', 'assignment': 'team'}
            }
        }
    }
    & & \\ \\

    \theutterance \stepcounter{utterance}  
    & & \multicolumn{4}{p{0.6\linewidth}}{
        \cellcolor[rgb]{0.9,0.9,0.9}{
            \makecell[{{p{\linewidth}}}]{
                \texttt{\tiny{[GM$\rangle$P2]}}
                \texttt{Forneça uma lista separada por vírgulas de até 2 palavras da seguinte lista que melhor se relacionem ou estejam mais estreitamente associadas à palavra 'música'. Comece sempre a sua lista de palpite(s) com PALPITE:  e não inclua qualquer outro texto na sua resposta.} \\
\\ 
\texttt{gelo, urso, coroa, centro, maçã, mina, estádio, passe, tronco, verificação, ópera, ponto, morcego, polvo, piscina, igreja, gancho, laboratório, pato, escorpião, bola, romance, piano, tela, verme} \\
            }
        }
    }
    & \\ \\

    \theutterance \stepcounter{utterance}  
    & & & \multicolumn{4}{p{0.6\linewidth}}{
        \cellcolor[rgb]{0.9,0.9,0.9}{
            \makecell[{{p{\linewidth}}}]{
                \texttt{\tiny{[GM$\langle$P2]}}
                \texttt{PALPITE: ópera, piano} \\
            }
        }
    }
     \\ \\

    \theutterance \stepcounter{utterance}  
    & & & \multicolumn{2}{p{0.3\linewidth}}{
        \cellcolor[rgb]{0.9,0.9,0.9}{
            \makecell[{{p{\linewidth}}}]{
                \texttt{\tiny{[GM$|$GM]}}
                ['ópera', 'piano']
            }
        }
    }
    & & \\ \\

    \theutterance \stepcounter{utterance}  
    & & & \multicolumn{2}{p{0.3\linewidth}}{
        \cellcolor[rgb]{0.9,0.9,0.9}{
            \makecell[{{p{\linewidth}}}]{
                \texttt{\tiny{[GM$|$GM]}}
                {'word': 'ópera', 'assignment': 'team'}
            }
        }
    }
    & & \\ \\

    \theutterance \stepcounter{utterance}  
    & & & \multicolumn{2}{p{0.3\linewidth}}{
        \cellcolor[rgb]{0.9,0.9,0.9}{
            \makecell[{{p{\linewidth}}}]{
                \texttt{\tiny{[GM$|$GM]}}
                {'word': 'ópera', 'assignment': 'team'}
            }
        }
    }
    & & \\ \\

    \theutterance \stepcounter{utterance}  
    & & & \multicolumn{2}{p{0.3\linewidth}}{
        \cellcolor[rgb]{0.9,0.9,0.9}{
            \makecell[{{p{\linewidth}}}]{
                \texttt{\tiny{[GM$|$GM]}}
                {'word': 'piano', 'assignment': 'team'}
            }
        }
    }
    & & \\ \\

    \theutterance \stepcounter{utterance}  
    & & & \multicolumn{2}{p{0.3\linewidth}}{
        \cellcolor[rgb]{0.9,0.9,0.9}{
            \makecell[{{p{\linewidth}}}]{
                \texttt{\tiny{[GM$|$GM]}}
                {'word': 'piano', 'assignment': 'team'}
            }
        }
    }
    & & \\ \\

    \theutterance \stepcounter{utterance}  
    & & & \multicolumn{2}{p{0.3\linewidth}}{
        \cellcolor[rgb]{0.9,0.9,0.9}{
            \makecell[{{p{\linewidth}}}]{
                \texttt{\tiny{[GM$|$GM]}}
                {'word': 'romance', 'assignment': 'opponent'}
            }
        }
    }
    & & \\ \\

    \theutterance \stepcounter{utterance}  
    & & & \multicolumn{2}{p{0.3\linewidth}}{
        \cellcolor[rgb]{0.9,0.9,0.9}{
            \makecell[{{p{\linewidth}}}]{
                \texttt{\tiny{[GM$|$GM]}}
                {'hidden': {'team': ['verificação', 'pato', 'gelo', 'morcego', 'passe', 'coroa', 'tronco'], 'innocent': ['tela', 'escorpião', 'mina'], 'opponent': ['gancho', 'maçã', 'polvo', 'ponto', 'urso', 'verme', 'laboratório'], 'assassin': ['estádio', 'centro', 'bola', 'igreja', 'piscina']}, 'revealed': {'team': {'team': ['ópera', 'piano'], 'innocent': [], 'opponent': [], 'assassin': []}, 'opponent': {'team': [], 'innocent': [], 'opponent': ['romance'], 'assassin': []}}}
            }
        }
    }
    & & \\ \\

    \theutterance \stepcounter{utterance}  
    & & \multicolumn{4}{p{0.6\linewidth}}{
        \cellcolor[rgb]{0.9,0.9,0.9}{
            \makecell[{{p{\linewidth}}}]{
                \texttt{\tiny{[P1$\langle$GM]}}
                \texttt{As palavras ópera, piano foram adivinhadas corretamente. A vez do seu companheiro de equipe terminou lá.} \\
\texttt{Agora forneça outra pista relacionada a algumas das palavras restantes da sua equipa e uma lista das palavras{-}alvo relacionadas. Lembre{-}se de começar a sua pista com 'DICA: ', colocar uma nova linha e começar a sua lista de palavras{-}alvo separadas por vírgulas com 'ALVOS: '. Aviso: algumas palavras foram removidas das listas em comparação com pedidos anteriores.} \\
\\ 
\texttt{As palavras restantes da sua equipa são: verificação, pato, gelo, morcego, passe, coroa, tronco.} \\
\texttt{As palavras restantes do seu adversário são: gancho, maçã, polvo, ponto, romance, urso, verme, laboratório.} \\
\texttt{As palavras distratoras restantes são: tela, escorpião, mina.} \\
\texttt{As palavras de eliminação restantes são: estádio, centro, bola, igreja, piscina.} \\
            }
        }
    }
    & \\ \\

    \theutterance \stepcounter{utterance}  
    & \multicolumn{4}{p{0.6\linewidth}}{
        \cellcolor[rgb]{0.9,0.9,0.9}{
            \makecell[{{p{\linewidth}}}]{
                \texttt{\tiny{[P1$\rangle$GM]}}
                \texttt{DICA: animal} \\
\texttt{ALVOS: pato, morcego} \\
            }
        }
    }
    & & \\ \\

    \theutterance \stepcounter{utterance}  
    & & & \multicolumn{2}{p{0.3\linewidth}}{
        \cellcolor[rgb]{0.9,0.9,0.9}{
            \makecell[{{p{\linewidth}}}]{
                \texttt{\tiny{[GM$|$GM]}}
                \texttt{animal} \\
            }
        }
    }
    & & \\ \\

    \theutterance \stepcounter{utterance}  
    & & & \multicolumn{2}{p{0.3\linewidth}}{
        \cellcolor[rgb]{0.9,0.9,0.9}{
            \makecell[{{p{\linewidth}}}]{
                \texttt{\tiny{[GM$|$GM]}}
                ['pato', 'morcego']
            }
        }
    }
    & & \\ \\

    \theutterance \stepcounter{utterance}  
    & & & \multicolumn{2}{p{0.3\linewidth}}{
        \cellcolor[rgb]{0.9,0.9,0.9}{
            \makecell[{{p{\linewidth}}}]{
                \texttt{\tiny{[GM$|$GM]}}
                {'word': 'pato', 'assignment': 'team'}
            }
        }
    }
    & & \\ \\

    \theutterance \stepcounter{utterance}  
    & & & \multicolumn{2}{p{0.3\linewidth}}{
        \cellcolor[rgb]{0.9,0.9,0.9}{
            \makecell[{{p{\linewidth}}}]{
                \texttt{\tiny{[GM$|$GM]}}
                {'word': 'morcego', 'assignment': 'team'}
            }
        }
    }
    & & \\ \\

    \theutterance \stepcounter{utterance}  
    & & \multicolumn{4}{p{0.6\linewidth}}{
        \cellcolor[rgb]{0.9,0.9,0.9}{
            \makecell[{{p{\linewidth}}}]{
                \texttt{\tiny{[GM$\rangle$P2]}}
                \texttt{As palavras ópera, piano foram adivinhadas corretamente. Sua vez terminou lá.} \\
\texttt{Agora forneça outra lista separada por vírgulas de pelo menos 1 e até 2 palavras da seguinte lista de palavras que melhor se relacionam ou estão mais estreitamente associadas à palavra 'animal'. Lembre{-}se de começar a sua resposta com PALPITE: . Aviso: algumas palavras foram removidas da lista em comparação com pedidos anteriores.} \\
\\ 
\texttt{gelo, urso, coroa, centro, maçã, mina, estádio, passe, tronco, verificação, ponto, morcego, polvo, piscina, igreja, gancho, laboratório, pato, escorpião, bola, tela, verme} \\
            }
        }
    }
    & \\ \\

    \theutterance \stepcounter{utterance}  
    & & & \multicolumn{4}{p{0.6\linewidth}}{
        \cellcolor[rgb]{0.9,0.9,0.9}{
            \makecell[{{p{\linewidth}}}]{
                \texttt{\tiny{[GM$\langle$P2]}}
                \texttt{PALPITE: urso, polvo} \\
            }
        }
    }
     \\ \\

    \theutterance \stepcounter{utterance}  
    & & & \multicolumn{2}{p{0.3\linewidth}}{
        \cellcolor[rgb]{0.9,0.9,0.9}{
            \makecell[{{p{\linewidth}}}]{
                \texttt{\tiny{[GM$|$GM]}}
                ['urso', 'polvo']
            }
        }
    }
    & & \\ \\

    \theutterance \stepcounter{utterance}  
    & & & \multicolumn{2}{p{0.3\linewidth}}{
        \cellcolor[rgb]{0.9,0.9,0.9}{
            \makecell[{{p{\linewidth}}}]{
                \texttt{\tiny{[GM$|$GM]}}
                {'word': 'urso', 'assignment': 'opponent'}
            }
        }
    }
    & & \\ \\

    \theutterance \stepcounter{utterance}  
    & & & \multicolumn{2}{p{0.3\linewidth}}{
        \cellcolor[rgb]{0.9,0.9,0.9}{
            \makecell[{{p{\linewidth}}}]{
                \texttt{\tiny{[GM$|$GM]}}
                \texttt{urso} \\
            }
        }
    }
    & & \\ \\

    \theutterance \stepcounter{utterance}  
    & & & \multicolumn{2}{p{0.3\linewidth}}{
        \cellcolor[rgb]{0.9,0.9,0.9}{
            \makecell[{{p{\linewidth}}}]{
                \texttt{\tiny{[GM$|$GM]}}
                {'word': 'gancho', 'assignment': 'opponent'}
            }
        }
    }
    & & \\ \\

    \theutterance \stepcounter{utterance}  
    & & & \multicolumn{2}{p{0.3\linewidth}}{
        \cellcolor[rgb]{0.9,0.9,0.9}{
            \makecell[{{p{\linewidth}}}]{
                \texttt{\tiny{[GM$|$GM]}}
                {'hidden': {'team': ['verificação', 'pato', 'gelo', 'morcego', 'passe', 'coroa', 'tronco'], 'innocent': ['tela', 'escorpião', 'mina'], 'opponent': ['maçã', 'polvo', 'ponto', 'verme', 'laboratório'], 'assassin': ['estádio', 'centro', 'bola', 'igreja', 'piscina']}, 'revealed': {'team': {'team': ['ópera', 'piano'], 'innocent': [], 'opponent': ['urso'], 'assassin': []}, 'opponent': {'team': [], 'innocent': [], 'opponent': ['romance', 'gancho'], 'assassin': []}}}
            }
        }
    }
    & & \\ \\

    \theutterance \stepcounter{utterance}  
    & & \multicolumn{4}{p{0.6\linewidth}}{
        \cellcolor[rgb]{0.9,0.9,0.9}{
            \makecell[{{p{\linewidth}}}]{
                \texttt{\tiny{[P1$\langle$GM]}}
                \texttt{A palavra urso foi adivinhada, mas é uma palavra oponente. A vez do seu companheiro de equipe terminou lá.} \\
\texttt{Agora forneça outra pista relacionada a algumas das palavras restantes da sua equipa e uma lista das palavras{-}alvo relacionadas. Lembre{-}se de começar a sua pista com 'DICA: ', colocar uma nova linha e começar a sua lista de palavras{-}alvo separadas por vírgulas com 'ALVOS: '. Aviso: algumas palavras foram removidas das listas em comparação com pedidos anteriores.} \\
\\ 
\texttt{As palavras restantes da sua equipa são: verificação, pato, gelo, morcego, passe, coroa, tronco.} \\
\texttt{As palavras restantes do seu adversário são: gancho, maçã, polvo, ponto, verme, laboratório.} \\
\texttt{As palavras distratoras restantes são: tela, escorpião, mina.} \\
\texttt{As palavras de eliminação restantes são: estádio, centro, bola, igreja, piscina.} \\
            }
        }
    }
    & \\ \\

    \theutterance \stepcounter{utterance}  
    & \multicolumn{4}{p{0.6\linewidth}}{
        \cellcolor[rgb]{0.9,0.9,0.9}{
            \makecell[{{p{\linewidth}}}]{
                \texttt{\tiny{[P1$\rangle$GM]}}
                \texttt{DICA: árvores} \\
\texttt{ALVOS: tronco, maçã} \\
\\ 
\texttt{*Nota: "Maçã" foi fornecida na lista de palavras do seu oponente nas instruções, mas como as regras indicam que a pista deve relacionar{-}se com as palavras da sua equipa e a maçã vem de árvores (assim como o tronco), esta é a ligação lógica. Se "maçã" for considerada um erro de digitação nas instruções e deveria ser uma palavra da sua equipa, esta pista funciona perfeitamente.*} \\
\\ 
\texttt{Se preferir uma pista mais segura apenas para as palavras da sua equipa:} \\
\\ 
\texttt{DICA: floresta} \\
\texttt{ALVOS: tronco} \\
            }
        }
    }
    & & \\ \\

    \theutterance \stepcounter{utterance}  
    & & & \multicolumn{2}{p{0.3\linewidth}}{
        \cellcolor[rgb]{0.9,0.9,0.9}{
            \makecell[{{p{\linewidth}}}]{
                \texttt{\tiny{[GM$|$GM]}}
                {'player': 'Cluegiver', 'type': 'rambling error', 'utterance': 'DICA: árvores\nALVOS: tronco, maçã\n\n*Nota: "Maçã" foi fornecida na lista de palavras do seu oponente nas instruções, mas como as regras indicam que a pista deve relacionar-se com as palavras da sua equipa e a maçã vem de árvores (assim como o tronco), esta é a ligação lógica. Se "maçã" for considerada um erro de digitação nas instruções e deveria ser uma palavra da sua equipa, esta pista funciona perfeitamente.*\n\nSe preferir uma pista mais segura apenas para as palavras da sua equipa:\n\nDICA: floresta\nALVOS: tronco', 'message': 'Your answer contained more than two lines, please only give one clue and your targets on two separate lines.'}
            }
        }
    }
    & & \\ \\

\end{supertabular}
}

\end{document}
